% ==============================================================================
% CEBE Interdisciplinary Fellowship Programme -- Call for PhD and PostDoc projects
% Project proposal template
%
% Formatting requirement from the call: Times New Roman, 12 pt,
% single-spaced text, 2.5 cm margins.
%
% Font choice:
%   This template uses `mathptmx`, a Times-metric-compatible font that
%   ships with every TeX distribution (no extra font files needed, works
%   the same on Windows/MiKTeX, Overleaf, and Linux TeX Live). It is
%   visually indistinguishable from Times New Roman and satisfies "Times
%   New Roman, 12pt" in practice.
%
%   If you would rather use the literal Times New Roman TTF installed on
%   this Windows machine, compile with XeLaTeX or LuaLaTeX instead and
%   swap the font block below for:
%       \usepackage{fontspec}
%       \setmainfont{Times New Roman}
%   (comment out \usepackage{mathptmx} in that case).
%
% Build: pdflatex -> bibtex -> pdflatex -> pdflatex   (or just `latexmk -pdf main.tex`)
% ==============================================================================

\documentclass[12pt,a4paper]{article}

% ---- Page geometry: 2.5 cm margins on all sides ----
\usepackage[a4paper,margin=2.5cm]{geometry}

% ---- Font: Times-compatible, 12pt (set via documentclass option above) ----
\usepackage[T1]{fontenc}
\usepackage[utf8]{inputenc}
\usepackage{mathptmx}

% ---- Single-spaced body text (this is already LaTeX's default, set explicitly) ----
\usepackage{setspace}
\singlespacing

% ---- Useful packages ----
\usepackage{graphicx}      % \includegraphics for a CEBE/university logo on the frontpage
\usepackage{enumitem}      % tighter, controllable lists
\usepackage{booktabs}      % clean tables (applicant list)
\usepackage[numbers,sort&compress]{natbib}  % references; bibliography style set below
\usepackage[colorlinks=true,linkcolor=blue,citecolor=blue,urlcolor=blue]{hyperref} % load near-last

% ---- Section numbering to mirror the call's own numbering (1, 2, 2.1 ... , 3, 4) ----
\setcounter{secnumdepth}{2}
\setcounter{tocdepth}{2}

% ---- Lightweight page-limit tracker -----------------------------------------
% Prints "PAGEMARK: <label> -- page N" to the compile log at each call. Grep
% the log (main.log) for "PAGEMARK" after compiling to check where each
% page-limited section starts/ends against the call's stated maxima:
%   Project Summary               max 1/2 page
%   Project Description           max 4 pages (excluding references)
%   Project relevance to CEBE     max 1/2 page
%   Sustainability goals          max 1/2 page
\newcommand{\pagemark}[1]{\typeout{PAGEMARK: #1 -- page \thepage}}

\begin{document}

% ==============================================================================
% FRONTPAGE
% Required content: title of project, name and affiliation of applicants,
% table of contents.
% ==============================================================================
\begin{titlepage}
\centering
\vspace*{1.5cm}

% Uncomment and point at a logo file if you want the CEBE / university mark here:
% \includegraphics[width=4cm]{cebe-logo.png}\par\vspace{1cm}

{\large CEBE Interdisciplinary Fellowship Programme}\par
{\large Call for PhD and PostDoc Projects}\par
\vspace{1.5cm}

{\Huge\bfseries [PROJECT TITLE]}\par
\vspace{0.5cm}
{\large [PhD project / Postdoc project -- delete as appropriate]}\par

\vspace{2cm}

\begin{tabular}{@{}p{3.2cm}p{4.5cm}p{6.5cm}@{}}
\toprule
\textbf{Role} & \textbf{Name} & \textbf{Affiliation} \\
\midrule
Principal Investigator (main supervisor/mentor) & [Name] & [Department, CEBE partner university] \\
Co-supervisor & [Name] & [Department, university] \\
Co-supervisor & [Name, if applicable] & [Department, university] \\
Candidate (if named) & [Name or ``to be recruited''] & [Affiliation] \\
\bottomrule
\end{tabular}

\vspace{2cm}
[Submission date]

\end{titlepage}

\newpage
\tableofcontents
\newpage

% ==============================================================================
% 1. PROJECT SUMMARY  (max 1/2 page)
% ==============================================================================
\pagemark{Project Summary -- start}
\section{Project Summary}

[Half a page, no more. One paragraph that a non-specialist reviewer can
read in two minutes: the problem, the interdisciplinary angle, the
approach, and the expected contribution.]

\pagemark{Project Summary -- end}

% ==============================================================================
% 2. PROJECT DESCRIPTION  (max 4 pages, excluding references)
% Required subsection headings -- keep exactly as below.
% ==============================================================================
\pagemark{Project Description -- start}
\section{Project Description}

\subsection{Motivation, Significance and Scientific Challenges}

[Clear description of the vision and goals, the distinguishing features
and foci of the project.]

[Motivation for why a doctoral student or postdoc specifically -- rather
than, e.g., a staff research project -- is the right vehicle for this
interdisciplinary work.]

\subsection{State of the Art}

[Current knowledge in each of the disciplines being bridged, and --
importantly -- the gap that only becomes visible/addressable at their
intersection.]

\subsection{Scientific Approach, Methodology, and Novelty}

[Research questions/hypotheses, methodology, work plan. Describe the
research contribution and its potential impact.]

\subsection{Research Environment and Supervision}

[Description of the research environment and infrastructure at each
participating group, demonstrating feasibility of the proposed project.]

[Description of the potential benefits and synergies expected through
the interdisciplinary collaboration -- be concrete: what does each side
bring that the other lacks?]

\subsection{Stakeholders involved in the project}

[Which industry partners and/or other stakeholders are involved, and
how -- e.g. co-supervision, data access, case studies, advisory role,
in-kind contribution, planned dissemination.]

\pagemark{Project Description -- end}

% References for the Project Description (kept here since the call's
% 4-page limit for this section explicitly excludes the reference list).
\bibliographystyle{unsrtnat}
\bibliography{references}

% ==============================================================================
% 3. PROJECT RELEVANCE TO CEBE RESEARCH FIELDS  (max 1/2 page)
% ==============================================================================
\pagemark{CEBE relevance -- start}
\section{Project relevance to CEBE research fields}

[Identify the relevant CEBE research field(s) from the list of seven
(see www.cebe.dk), and describe the expected contribution to that
field / those fields. Be explicit and specific rather than mapping to
all seven.]

\pagemark{CEBE relevance -- end}

% ==============================================================================
% 4. SUSTAINABILITY GOALS AND RELATION TO THE CEBE PROGRAM VISION  (max 1/2 page)
% ==============================================================================
\pagemark{Sustainability goals -- start}
\section{Sustainability goals of the project and relation to the CEBE program vision}

[Specific sustainability objectives, including any sustainability-related
trade-offs the project may involve.]

[Sustainability-related conceptual and methodological choices made in
the project design.]

\pagemark{Sustainability goals -- end}

\end{document}
